\section{P1：二维候选路径生成}

P1 面向给定二维仓库场景，为关键节点对生成可行候选路径及其成本、几何记录和轨迹采样表。其输入包括节点、障碍物、功能区域、AMR 初始点和任务取送货点；输出包括路径成本表、成本矩阵、路径几何记录和相对时间轨迹采样表。P1 的建模重点是把连续仓库平面中的可达性、障碍物膨胀、转向偏好和动态风险代价统一写入可复现的路径生成结果。

\subsection{建模对象与输出目标}

P1 的输入来自 P0 整理后的二维仓库场景，主要包括节点表、区域表、障碍物表、AMR 初始位置和任务取送货点。关键节点集合记为 $V_K$，其中包含入库口、出库口、货架点、分拣台、充电点和 AMR 初始点。P1 对所有有序节点对 $(i,j), i\neq j$ 生成路径，因此每一种单算法在本算例中输出 $17\times 16=272$ 条可行路径。

从数学上看，P1 需要为每个 OD 对输出一条折线路径
\[
  p_{ij} = (q_0,q_1,\ldots,q_m),\qquad q_0=i,\ q_m=j,
\]
其中每个 $q_\ell=(x_\ell,y_\ell)$ 是二维平面上的路径点。路径的几何距离为
\[
  d(p_{ij})=\sum_{\ell=0}^{m-1}\|q_{\ell+1}-q_\ell\|_2,
\]
在本文实现中 AMR 标准速度取 $v=1.0$，所以物理通行时间为
\[
  t(p_{ij})=\frac{d(p_{ij})}{v}.
\]
除了距离和时间，P1 还记录转角数、转角代价、动态区域代价、轨迹采样点数等指标。这些指标一方面用于算法比较，另一方面用于路径选择、轨迹展开和可复现实验检查。

障碍物处理采用“AMR footprint 膨胀”思想。货架、墙体和动态封锁区在路径规划时不是只按原始几何边界处理，而是按机器人半径进行膨胀，使路径几何中心线与障碍物保持安全距离。这样，路径展开为 footprint 轨迹时不会天然穿越障碍区域。

\subsection{算法框架}

P1 实现了 \verb|basic_astar|、\verb|vg|、\verb|avg|、\verb|davg| 四类路径生成算法，分别代表不同路径偏好：栅格可达性、几何短路、转弯平滑性和动态风险规避。P1 在本节只负责生成这些单算法路径源；多算法融合得到 mixed 多候选路径库的工作属于 P2 的调度预处理。

\begin{table}[H]
\centering
\small
\begin{tabular}{p{0.18\linewidth}p{0.34\linewidth}p{0.36\linewidth}}
\toprule
算法 & 核心思想 & 当前实现中的作用 \\
\midrule
\verb|basic_astar| & 将仓库离散为 $0.1$m 栅格，在 8 邻域上用 A* 最小化几何距离 & 作为最朴素的可达性基线，路径贴合栅格方向，通常折点较多 \\
\verb|vg| & 构造障碍物膨胀边界的可视图，在可见顶点之间搜索最短折线 & 作为连续几何最短路基线，路径通常更短、更少折点 \\
\verb|avg| & 在 VG 基础上把上一顶点纳入搜索状态，对转角大小加惩罚 & 倾向选择转向更少、更平滑的路径，体现运动执行偏好 \\
\verb|davg| & 在 AVG 基础上加入动态风险区域暴露代价，并采用局部活跃可视图 & 用于动态封锁或风险场景，倾向绕开临时风险区域 \\

\bottomrule
\end{tabular}
\caption{P1 路径算法在最终系统中的定位}
\label{tab:p1-methods}
\end{table}

表 \ref{tab:p1-methods} 说明 P1 并不是简单比较几种最短路算法，而是生成不同偏好的基础路径源。A* 偏重栅格可达性，VG 偏重连续空间中的几何最短，AVG 引入转弯平滑性，DAVG 引入动态风险代价。这些单算法结果可以进一步融合、去重并构造 mixed 多候选路径库。

\subsection{Grid A*：栅格可达性基线}

\verb|basic_astar| 将二维仓库转化为分辨率 $0.1$m 的规则栅格。每个栅格单元用其中心点代表空间位置，若中心点落入膨胀后的硬障碍物，则该单元标记为不可通行。搜索动作采用 8 邻域：
\[
  \Delta=\{(-1,0),(1,0),(0,-1),(0,1),(-1,-1),(-1,1),(1,-1),(1,1)\}.
\]
横纵向移动代价为一个栅格宽度，斜向移动代价为 $\sqrt{2}$ 个栅格宽度。A* 的评价函数为
\[
  f(n)=g(n)+h(n),
\]
其中 $g(n)$ 是从起点到当前栅格的累计距离，$h(n)$ 是当前栅格到终点的欧氏直线距离。由于 8 邻域移动的真实最短剩余距离不会小于欧氏距离，这个启发函数是可接受的，不会高估剩余代价。

实现中还加入了两个工程细节。第一，如果起点或终点恰好位于膨胀障碍物内，会在局部半径内寻找最近可通行栅格，避免由于节点坐标贴近货架边界而直接失败。第二，斜向移动时检查相邻两个正交单元，禁止“贴角穿越”，防止路径从两个障碍物角点之间不合理穿过。

Grid A* 的优点是直观、稳定、容易处理任意形状的占用区域；缺点是路径受栅格方向限制，可能产生较多方向变化。本文只把转角数作为报告指标，不把转角惩罚加入 A* 搜索目标，因此它是一个纯距离基线。后续的 VG/AVG/DAVG 都可以与该基线对比，说明连续几何建图和路径偏好项的价值。

\subsection{VG：可视图几何最短路}

\verb|vg| 不再在密集栅格上扩展，而是在连续平面中构造稀疏可视图。可视图的顶点包括起点、终点和膨胀障碍物的角点。若两个顶点之间的线段不穿过任何硬障碍物内部，则二者之间连一条边，边权为欧氏距离。随后在这个可视图上运行 A*，启发函数仍是到终点的直线距离。

VG 的关键是把“绕障碍物”问题转化为“在可见角点之间选折线”的图搜索问题。若不考虑转弯和动态风险，二维多边形障碍物环境中的最短折线通常会贴近障碍物膨胀边界的顶点。因此，VG 比 Grid A* 更接近连续空间中的几何最短路，路径折点也少得多。

本文实现中，线段可见性使用矩形裁剪思想判断线段是否穿过障碍物内部。触碰边界被允许，因为障碍物已经经过安全膨胀；真正被拒绝的是线段与膨胀障碍物内部有正长度重叠的情况。墙体作为仓库边界，不参与候选绕行角点，避免路径沿外墙边界产生无意义顶点。

VG 的输出中 \verb|turn_cost|、\verb|dynamic_cost| 均为 0，\verb|total_cost| 等于几何距离。它在 P1 里提供的是“最短可见折线”的参照物：如果某条路径在 VG 中已经很短且折点很少，说明该 OD 对的空间约束比较直接；如果与 A* 差异很大，则说明栅格离散或障碍边界对路径形态影响明显。

\subsection{AVG：加入转角偏好的增强可视图}

普通 VG 只关心路径长度，可能选择一条距离略短但转弯较急的路线。对于 AMR 来说，频繁或剧烈转向会带来控制代价、速度损失和执行不稳定性。AVG 在 VG 的几何图基础上引入转角惩罚，目标函数变为
\[
  c_{\mathrm{AVG}}(p)=d(p)+\lambda_\theta\sum_{\ell=1}^{m-1}\theta_\ell,
\]
其中 $\theta_\ell$ 是路径在顶点 $q_\ell$ 的航向变化角，$\lambda_\theta=0.5$ 为单位弧度转角代价。

为了在搜索过程中正确计算转角，AVG 的状态不再只是当前顶点，而是
\[
  s=(q_{\mathrm{prev}},q_{\mathrm{cur}}).
\]
当搜索从 $q_{\mathrm{cur}}$ 扩展到 $q_{\mathrm{next}}$ 时，就可以根据上一段和下一段的航向计算当前顶点处的转角代价。该状态扩展使 A* 的转移代价从单纯边长变为
\[
  c(s,q_{\mathrm{next}})=\|q_{\mathrm{next}}-q_{\mathrm{cur}}\|_2
  +\lambda_\theta\,\Delta\phi(q_{\mathrm{prev}},q_{\mathrm{cur}},q_{\mathrm{next}}).
\]
第一步没有上一段航向，因此不收取转角代价；同时实现中跳过立即掉头的扩展，减少无意义循环。

需要强调的是，AVG 中的转角代价不是物理通行时间。本文仍用几何距离除以速度计算 \verb|travel_time|，转角项只进入 \verb|total_cost|，用于表达路径偏好。这样做的好处是保持时间窗模型简洁，同时让 P2 在路径选择时可以区分“距离短但转弯多”和“略长但更平顺”的路线。

\subsection{DAVG：动态风险代价与活跃可视图}

DAVG 在 AVG 的基础上处理动态封锁或风险区域。它保留状态增强和转角惩罚，同时对靠近动态区域的可视边加入暴露代价：
\[
  c_{\mathrm{DAVG}}(p)=d(p)+\lambda_\theta\sum_{\ell=1}^{m-1}\theta_\ell
  +\sum_{e\in p} c_{\mathrm{dyn}}(e).
\]
对于一条边 $e$，若它到动态区域的最小距离为 $\delta_e$，动态风险缓冲半径为 $B=1.0$，权重为 $\lambda_d=2.0$，则实现中的动态代价可理解为
\[
  c_{\mathrm{dyn}}(e)=
  \begin{cases}
    \lambda_d\left(1-\frac{\delta_e}{B}\right)|e|, & \delta_e<B,\\
    0, & \delta_e\ge B.
  \end{cases}
\]
这意味着边越长、越靠近动态区域，付出的风险代价越大。若路径远离动态区域，则 DAVG 会退化为 AVG。

DAVG 还加入活跃区域建图：对每个 OD 对，只把起终点走廊附近的障碍物角点加入候选顶点，同时动态障碍物角点始终保留。这样可以减少远离当前 OD 的障碍角点带来的图规模。为了不牺牲可行性，可见性检测仍对所有硬障碍物进行；若局部活跃图无法找到路径，则回退到完整障碍角点图。

在当前静态主算例中，DAVG 的平均动态代价为 0，说明选出的路径没有进入动态风险缓冲区。因此，DAVG 与 AVG 的几何距离、转角数和总成本在 P1 层面基本一致。这不是 DAVG 失效，而是说明当前静态路径比较没有触发风险绕行；DAVG 的风险绕行价值会在动态事件或封锁场景中体现。

\subsection{输出接口}

P1 的主要输出位于 \verb|data/processed/p1/<algorithm>/|：

\begin{itemize}
  \item \verb|key_nodes.csv|：参与路径规划的关键节点；
  \item \verb|path_cost.csv|：OD 路径成本表；
  \item \verb|path_grid_cells.csv|：路径穿越栅格或可视图顶点记录；
  \item \verb|path_trajectory_samples.csv|：相对时间轨迹采样；
  \item \verb|path_cost_matrix_time.csv| 和 \verb|path_cost_matrix_total.csv|：时间矩阵和综合成本矩阵。
\end{itemize}

其中，\verb|path_cost.csv| 提供 \verb|from_node|、\verb|to_node|、\verb|path_uid|、\verb|distance|、\verb|travel_time|、\verb|turn_cost|、\verb|dynamic_cost|、\verb|total_cost| 和 \verb|planning_status|。其中 \verb|planning_status| 用于标记 OD 对是否成功生成路径，避免不可达路径被误用。

\verb|path_trajectory_samples.csv| 提供相对出发时刻的二维轨迹样本，包括 \verb|offset_time|、\verb|x|、\verb|y|、\verb|theta|、\verb|speed| 和 \verb|footprint_radius|。该表保留路径上的几何点、朝向、速度和机器人 footprint 半径，便于后续按具体任务开始时间展开为绝对时空轨迹。

通过上述表结构，P1 完成了从“二维空间路径”到“路径成本矩阵”和“轨迹采样数据”的转换。路径几何用于解释算法差异，时间矩阵用于记录 OD 对通行时间，轨迹采样用于保留路径在二维平面中的可执行形态。
